\FigureLayoutDeclare{img/1990/national_science/q5_fig1.png}{4.2cm}{68bp 34bp 175bp 53bp}

\chapter{1990年全国卷（理）}
\section{单选题}

\begin{problem}
方程 \(2^{\log_3 x}=\frac{1}{4}\) 的解是（\quad）

\choices
    {\(x=\frac{1}{9}\)}
    {\(x=\frac{\sqrt{3}}{3}\)}
    {\(x=\sqrt{3}\)}
    {\(x=9\)}

\begin{answer}
A
\end{answer}

\begin{solution}
因 \(2^t\) 在 \(\R\) 上严格递增，令 \(t=\log_3 x\)，由 \(2^t=\frac{1}{4}=2^{-2}\) 得 \(t=-2\)，所以 \(\log_3 x=-2\)，从而 \(x=3^{-2}=\frac{1}{9}\). 该值满足 \(x>0\)，故选 A.
\end{solution}
\end{problem}

\begin{problem}
把复数 \(1+\i\) 对应的向量按顺时针方向旋转 \(\frac{2\pi}{3}\)，所得到的向量对应的复数是（\quad）

\choices
    {\(\frac{1-\sqrt{3}}{2}+\frac{-1+\sqrt{3}}{2}\i\)}
    {\(\frac{-1+\sqrt{3}}{2}+\frac{-1-\sqrt{3}}{2}\i\)}
    {\(\frac{-1+\sqrt{3}}{2}+\frac{1-\sqrt{3}}{2}\i\)}
    {\(\frac{1-\sqrt{3}}{2}+\frac{-1-\sqrt{3}}{2}\i\)}

\begin{answer}
B
\end{answer}

\begin{solution}
顺时针旋转 \(\frac{2\pi}{3}\) 相当于乘以 \(\cos\frac{2\pi}{3}-\i\sin\frac{2\pi}{3}=-\frac{1}{2}-\frac{\sqrt{3}}{2}\i\). 因而所得复数为 \((1+\i)\left(-\frac{1}{2}-\frac{\sqrt{3}}{2}\i\right)=\frac{\sqrt{3}-1}{2}-\frac{1+\sqrt{3}}{2}\i\)，故选 B.
\end{solution}
\end{problem}

\begin{problem}
如果轴截面为正方形的圆柱的侧面积是 \(S\)，那么圆柱的体积等于（\quad）

\choices
    {\(\frac{S}{2}\sqrt{S}\)}
    {\(\frac{S}{2}\sqrt{\frac{S}{\pi}}\)}
    {\(\frac{S}{4}\sqrt{S}\)}
    {\(\frac{S}{4}\sqrt{\frac{S}{\pi}}\)}

\begin{answer}
D
\end{answer}

\begin{solution}
设圆柱底面半径为 \(r\)，高为 \(h\). 轴截面是边长分别为 \(2r\) 与 \(h\) 的矩形，由轴截面为正方形得 \(h=2r\). 因而侧面积 \(S=2\pi rh=4\pi r^2\)，所以 \(r=\frac{1}{2}\sqrt{\frac{S}{\pi}}\). 圆柱体积为 \(V=\pi r^2h=2\pi r^3=\frac{S}{4}\sqrt{\frac{S}{\pi}}\)，故选 D.
\end{solution}
\end{problem}

\begin{problem}
方程 \(\sin 2x=\sin x\) 在区间 \((0,2\pi)\) 内的解个数是（\quad）

\choices
    {\(1\)}
    {\(2\)}
    {\(3\)}
    {\(4\)}

\begin{answer}
C
\end{answer}

\begin{solution}
方程可化为 \(\sin 2x-\sin x=0\)，即 \(\sin x(2\cos x-1)=0\). 在区间 \((0,2\pi)\) 内，\(\sin x=0\) 给出 \(x=\pi\)，而 \(\cos x=\frac{1}{2}\) 给出 \(x=\frac{\pi}{3}\)，\(\frac{5\pi}{3}\). 这三个解互不相同，所以解的个数为 \(3\)，故选 C.
\end{solution}
\end{problem}

\begin{problem}
如图是函数 \(y=2\sin\left( \omega x+\varphi \right)\)（\(\left| \varphi \right|<\frac{\pi}{2}\)）的图象，那么（\quad）

\bitmapfigure[max width=0.34\linewidth]{img/1990/national_science/q5_fig1.png}

\choices
    {\(\omega=\frac{10}{11}\)，\(\varphi=\frac{\pi}{6}\)}
    {\(\omega=\frac{10}{11}\)，\(\varphi=-\frac{\pi}{6}\)}
    {\(\omega=2\)，\(\varphi=\frac{\pi}{6}\)}
    {\(\omega=2\)，\(\varphi=-\frac{\pi}{6}\)}

\begin{answer}
C
\end{answer}

\begin{solution}
由图象在 \(x=0\) 处的函数值为 \(1\)，得 \(2\sin\varphi=1\). 结合 \(\lvert\varphi\rvert<\frac{\pi}{2}\)，可知 \(\varphi=\frac{\pi}{6}\). 图象还显示在 \(x=\frac{11\pi}{12}\) 处曲线由下向上经过 \(x\) 轴，所以此时相位为 \(2\pi\)，从而 \(\omega\frac{11\pi}{12}+\frac{\pi}{6}=2\pi\)，解得 \(\omega=2\). 因此选 C.
\end{solution}
\end{problem}

\begin{problem}
函数 \(y=\frac{\sin x}{\left| \sin x \right|}+\frac{\left| \cos x \right|}{\cos x}+\frac{\tan x}{\left| \tan x \right|}+\frac{\left| \cot x \right|}{\cot x}\) 的值域是（\quad）

\choices
    {\(\{-2,4\}\)}
    {\(\{-2,0,4\}\)}
    {\(\{-2,0,2,4\}\)}
    {\(\{-4,-2,0,4\}\)}

\begin{answer}
B
\end{answer}

\begin{solution}
定义域要求 \(\sin x\neq0\) 且 \(\cos x\neq0\). 记 \(s=\frac{\sin x}{\lvert\sin x\rvert}\)，\(c=\frac{\lvert\cos x\rvert}{\cos x}\)，则 \(s\)，\(c\in\{1,-1\}\)，且 \(\frac{\tan x}{\lvert\tan x\rvert}=sc\)，\(\frac{\lvert\cot x\rvert}{\cot x}=sc\). 所以函数值为 \(s+c+2sc\). 在第一、二、三、四象限中，\((s,c)\) 分别为 \((1,1)\)、\((1,-1)\)、\((-1,-1)\)、\((-1,1)\)，对应函数值分别为 \(4\)、\(-2\)、\(0\)、\(-2\). 因此值域为 \(\{-2,0,4\}\)，故选 B.
\end{solution}
\end{problem}

\begin{problem}
如果直线 \(y=ax+2\) 与直线 \(y=3x-b\) 关于直线 \(y=x\) 对称，那么（\quad）

\choices
    {\(a=\frac{1}{3}\)，\(b=6\)}
    {\(a=\frac{1}{3}\)，\(b=-6\)}
    {\(a=3\)，\(b=-2\)}
    {\(a=3\)，\(b=6\)}

\begin{answer}
A
\end{answer}

\begin{solution}
关于直线 \(y=x\) 的对称变换交换点的两个坐标. 设第一条直线上的点 \((x,y)\) 变为 \((X,Y)=(y,x)\)，则 \(X=aY+2\)，即 \(Y=\frac{1}{a}X-\frac{2}{a}\). 该直线与 \(Y=3X-b\) 重合，因此 \(\frac{1}{a}=3\)，且 \(-\frac{2}{a}=-b\)，解得 \(a=\frac{1}{3}\)，\(b=6\)，故选 A.
\end{solution}
\end{problem}

\begin{problem}
极坐标方程 \(4\rho\sin^2\frac{\theta}{2}=5\) 表示的曲线是（\quad）

\choices
    {圆}
    {椭圆}
    {双曲线的一支}
    {抛物线}

\begin{answer}
D
\end{answer}

\begin{solution}
由 \(\sin^2\frac{\theta}{2}=\frac{1-\cos\theta}{2}\)，原方程化为 \(2\rho(1-\cos\theta)=5\)，即 \(\rho(1-\cos\theta)=\frac{5}{2}\). 由于 \(x=\rho\cos\theta\)，所以 \(\rho-x=\frac{5}{2}\)，即 \(\rho=x+\frac{5}{2}\). 这里 \(\rho\) 是点到原点的距离，而 \(x+\frac{5}{2}\) 是点到直线 \(x=-\frac{5}{2}\) 的距离，因此该曲线是以原点为焦点、以 \(x=-\frac{5}{2}\) 为准线的抛物线，故选 D.
\end{solution}
\end{problem}

\begin{problem}
设全集 \(I=\left\{(x,y)\mid x,y\in\R\right\}\)，集合 \(M=\left\{(x,y)\mid \frac{y-3}{x-2}=1\right\}\)，\(N=\left\{(x,y)\mid y\neq x+1\right\}\)，那么 \(\overline{M}\cap \overline{N}\) 等于（\quad）

\choices
    {\(\varnothing\)}
    {\(\{(2,3)\}\)}
    {\((2,3)\)}
    {\(\{(x,y)\mid y=x+1\}\)}

\begin{answer}
B
\end{answer}

\begin{solution}
记 \(L=\{(x,y)\mid y=x+1\}\)，上划线表示相对于全集 \(I\) 的补集. 由 \(\frac{y-3}{x-2}=1\) 可知 \(x\neq2\)，且 \(y=x+1\)，所以 \(M=L\smallsetminus\{(2,3)\}\). 又 \(N=I\smallsetminus L\)，从而 \(\overline{N}=L\)，\(\overline{M}=(I\smallsetminus L)\cup\{(2,3)\}\). 因此 \(\overline{M}\cap\overline{N}=\{(2,3)\}\)，故选 B.
\end{solution}
\end{problem}

\begin{problem}
如果实数 \(x\)，\(y\) 满足等式 \((x-2)^2+y^2=3\)，那么 \(\frac{y}{x}\) 的最大值是（\quad）

\choices
    {\(\frac{1}{2}\)}
    {\(\frac{\sqrt{3}}{3}\)}
    {\(\frac{\sqrt{3}}{2}\)}
    {\(\sqrt{3}\)}

\begin{answer}
D
\end{answer}

\begin{solution}
圆的横坐标范围为 \([2-\sqrt{3},2+\sqrt{3}]\)，且 \(2-\sqrt{3}>0\)，所以圆上任一点的 \(x\) 坐标均为正. 令 \(m=\frac{y}{x}\)，即 \(y=mx\)，代入 \((x-2)^2+y^2=3\) 得 \((1+m^2)x^2-4x+1=0\). 要使 \(x\) 为实数，必须有判别式 \(16-4(1+m^2)\geqslant0\)，从而 \(m^2\leqslant3\)，所以 \(m\leqslant\sqrt{3}\). 当 \(m=\sqrt{3}\) 时判别式为零，取 \(x=\frac{1}{2}\)，\(y=\frac{\sqrt{3}}{2}\) 即满足圆的方程，故最大值为 \(\sqrt{3}\)，选 D.
\end{solution}
\end{problem}

\begin{problem}
如图，正三棱锥 \(S-ABC\) 的侧棱与底面边长相等，如果 \(E\)，\(F\) 分别为 \(SC\)，\(AB\) 的中点，那么异面直线 \(EF\) 与 \(SA\) 所成的角等于（\quad）

\bitmapfigure[max width=0.30\linewidth]{img/1990/national_science/q11_fig1.png}

\choices
    {\(90^\circ\)}
    {\(60^\circ\)}
    {\(45^\circ\)}
    {\(30^\circ\)}

\begin{answer}
C
\end{answer}

\begin{solution}
取棱长为 \(1\)，建立直角坐标系，令 \(A=(0,0,0)\)，\(B=(1,0,0)\)，\(C=(\frac{1}{2},\frac{\sqrt{3}}{2},0)\)，\(S=(\frac{1}{2},\frac{\sqrt{3}}{6},\sqrt{\frac{2}{3}})\). 这四点构成正四面体. 于是 \(E=(\frac{1}{2},\frac{\sqrt{3}}{3},\frac{1}{2}\sqrt{\frac{2}{3}})\)，\(F=(\frac{1}{2},0,0)\)，所以 \(\overrightarrow{FE}=(0,\frac{\sqrt{3}}{3},\frac{1}{2}\sqrt{\frac{2}{3}})\)，\(\overrightarrow{AS}=(\frac{1}{2},\frac{\sqrt{3}}{6},\sqrt{\frac{2}{3}})\). 计算得 \(\overrightarrow{FE}\cdot\overrightarrow{AS}=\frac{1}{2}\)，\(\lvert\overrightarrow{FE}\rvert^2=\frac{1}{2}\)，且 \(\lvert\overrightarrow{AS}\rvert=1\). 设异面直线 \(EF\) 与 \(SA\) 所成的锐角为 \(\theta\)，则 \(\cos\theta=\frac{\lvert\overrightarrow{FE}\cdot\overrightarrow{AS}\rvert}{\lvert\overrightarrow{FE}\rvert\lvert\overrightarrow{AS}\rvert}=\frac{1}{\sqrt{2}}\)，所以 \(\theta=45^\circ\)，故选 C.
\end{solution}
\end{problem}

\begin{problem}
已知 \(h>0\). 设命题甲为：两个实数 \(a\)，\(b\) 满足 \(|a-b|<2h\)；命题乙为：两个实数 \(a\)，\(b\) 满足 \(|a-1|<h\) 且 \(|b-1|<h\). 那么甲是乙的（\quad）

\choices
    {充分不必要条件}
    {必要不充分条件}
    {充要条件}
    {既不充分也不必要条件}

\begin{answer}
B
\end{answer}

\begin{solution}
若命题乙成立，则由三角不等式有 \(\lvert a-b\rvert\leqslant\lvert a-1\rvert+\lvert b-1\rvert<2h\)，所以乙可以推出甲，甲是乙成立的必要条件. 但甲不能推出乙，例如取 \(a=b=1+2h\)，有 \(\lvert a-b\rvert=0<2h\)，而 \(\lvert a-1\rvert=2h>h\)，乙不成立. 因此甲是乙的必要不充分条件，故选 B.
\end{solution}
\end{problem}

\begin{problem}
A，B，C，D，E 五人并排站成一排，如果 B 必须站在 A 的右边（A，B 可以不相邻），那么不同的排法共有（\quad）

\choices
    {\(24\text{ 种}\)}
    {\(60\text{ 种}\)}
    {\(90\text{ 种}\)}
    {\(120\text{ 种}\)}

\begin{answer}
B
\end{answer}

\begin{solution}
五个人的全部排法共有 \(5!=120\) 种. 对任意一种排法，交换 A 与 B 的位置可得到唯一的另一种排法，并且这两种排法中恰有一种满足 B 在 A 的右边. 因而满足条件的排法数为 \(\frac{5!}{2}=60\)，故选 B.
\end{solution}
\end{problem}

\begin{problem}
以一个正方体的顶点为顶点的四面体共有（\quad）

\choices
    {\(70\text{ 个}\)}
    {\(64\text{ 个}\)}
    {\(58\text{ 个}\)}
    {\(52\text{ 个}\)}

\begin{answer}
C
\end{answer}

\begin{solution}
把正方体的八个顶点设为 \((x,y,z)\)，其中 \(x\)，\(y\)，\(z\in\{0,1\}\). 任取四个顶点共有 \(\mathrm C_8^4=70\) 组. 其中共面的四点包括六个面各自的四个顶点；此外，非面共面的四点所在的平面分别为 \(x=y\)、\(x+y=1\)、\(x=z\)、\(x+z=1\)、\(y=z\)、\(y+z=1\)，共六组. 这些共面四点各对应一个退化的四面体，且上述十二组互不重复. 因此四面体的个数为 \(\mathrm C_8^4-6-6=70-12=58\)，故选 C.
\end{solution}
\end{problem}

\begin{problem}
设函数 \(y=\arctan x\) 的图象沿 \(x\) 轴正方向平移 \(2\) 个单位所得到的图象为 \(C\). 又设图象 \(C'\) 与 \(C\) 关于原点对称，那么 \(C'\) 所对应的函数是（\quad）

\choices
    {\(y=-\arctan(x-2)\)}
    {\(y=\arctan(x-2)\)}
    {\(y=-\arctan(x+2)\)}
    {\(y=\arctan(x+2)\)}

\begin{answer}
D
\end{answer}

\begin{solution}
图象 \(C\) 是 \(y=\arctan x\) 沿 \(x\) 轴正方向平移 \(2\) 个单位得到的，所以其函数为 \(y=\arctan(x-2)\). 若 \((X,Y)\) 在关于原点对称的图象 \(C'\) 上，则 \((-X,-Y)\) 在 \(C\) 上，故 \(-Y=\arctan(-X-2)\). 利用 \(\arctan\) 为奇函数，得 \(Y=\arctan(X+2)\)，所以 \(C'\) 所对应的函数为 \(y=\arctan(x+2)\)，故选 D.
\end{solution}
\end{problem}

\section{填空题}

\begin{problem}
双曲线 \(\frac{y^2}{16}-\frac{x^2}{9}=1\) 的准线方程是\fillinblank{}.
\end{problem}

\begin{answer}
\(y=\pm\frac{16}{5}\)
\end{answer}

\begin{solution}
将双曲线写成标准形式 \(\frac{y^2}{a^2}-\frac{x^2}{b^2}=1\)，可得 \(a=4\)，\(b=3\). 由 \(c^2=a^2+b^2\) 得 \(c=5\)，离心率为 \(e=\frac{c}{a}=\frac{5}{4}\). 双曲线的准线方程为 \(y=\pm\frac{a}{e}\)，所以 \(y=\pm\frac{4}{5/4}=\pm\frac{16}{5}\).
\end{solution}

\begin{problem}
\((x-1)-(x-1)^2+(x-1)^3-(x-1)^4+(x-1)^5\) 的展开式中，\(x^2\) 的系数等于\fillinblank{}.
\end{problem}

\begin{answer}
\(-20\)
\end{answer}

\begin{solution}
令 \(t=x-1\). 原式为 \(t-t^2+t^3-t^4+t^5\). 其中 \(t\) 不含 \(x^2\) 项；在 \(-t^2\)、\(t^3\)、\(-t^4\)、\(t^5\) 中，\(x^2\) 的系数依次为 \(-1\)、\(\mathrm C_3^2(-1)=-3\)、\(-\mathrm C_4^2=-6\)、\(\mathrm C_5^2(-1)^3=-10\). 所以所求系数为 \(-1-3-6-10=-20\).
\end{solution}

\begin{problem}
已知 \(\{a_n\}\) 是公差不为零的等差数列，如果 \(S_n\) 是 \(\{a_n\}\) 的前 \(n\) 项的和，那么 \(\displaystyle\lim_{n\to\infty}\frac{na_n}{S_n}=\)\fillinblank{}.
\end{problem}

\begin{answer}
\(2\)
\end{answer}

\begin{solution}
设等差数列首项为 \(a_1\)，公差为 \(d\)，其中 \(d\neq0\). 则 \(a_n=a_1+(n-1)d=dn+a_1-d\)，且 \(S_n=\frac{n}{2}(a_1+a_n)=\frac{n}{2}(dn+2a_1-d)\). 因为 \(d\neq0\)，一次式 \(dn+2a_1-d\) 至多有一个零点，所以当 \(n\) 足够大时分式恒有意义. 对这些项约去一个 \(n\)，得 \(\frac{na_n}{S_n}=\frac{2(dn+a_1-d)}{dn+2a_1-d}=\frac{2\left(d+\frac{a_1-d}{n}\right)}{d+\frac{2a_1-d}{n}}\). 令 \(n\to\infty\)，得到 \(\displaystyle\lim_{n\to\infty}\frac{na_n}{S_n}=2\).
\end{solution}

\begin{problem}
函数 \(y=\sin x\cos x+\sin x+\cos x\) 的最大值是\fillinblank{}.
\end{problem}

\begin{answer}
\(\frac{1}{2}+\sqrt{2}\)
\end{answer}

\begin{solution}
令 \(u=\sin x+\cos x\)，则 \(-\sqrt{2}\leqslant u\leqslant\sqrt{2}\)，并且 \(\sin x\cos x=\frac{u^2-1}{2}\). 所以原函数值为 \(\frac{u^2-1}{2}+u=\frac{u^2}{2}+u-\frac{1}{2}\). 这是开口向上的二次函数，在区间 \([-\sqrt{2},\sqrt{2}]\) 上的最大值取在端点；两端点的函数值分别为 \(\frac{1}{2}-\sqrt{2}\) 与 \(\frac{1}{2}+\sqrt{2}\). 因此最大值为 \(\frac{1}{2}+\sqrt{2}\).
\end{solution}

\begin{problem}
如图，三棱柱 \(ABC-A_1B_1C_1\) 中，若 \(E\)，\(F\) 分别为 \(AB\)，\(AC\) 的中点，平面 \(EB_1C_1F\) 将三棱柱分成体积为 \(V_1\)，\(V_2\) 的两部分，那么 \(V_1:V_2=\)\fillinblank{}.

\bitmapfigure[max width=0.24\linewidth]{img/1990/national_science/q20_fig1.png}
\end{problem}

\begin{answer}
\(7:5\)
\end{answer}

\begin{solution}
体积比在可逆仿射变换下不变，取三棱柱的坐标为 \(A=(0,0,0)\)，\(B=(1,0,0)\)，\(C=(0,1,0)\)，\(A_1=(0,0,1)\)，\(B_1=(1,0,1)\)，\(C_1=(0,1,1)\). 此时 \(E=(\frac{1}{2},0,0)\)，\(F=(0,\frac{1}{2},0)\). 平面 \(EB_1C_1F\) 上的四点都满足 \(2x+2y-z=1\)，故其方程为 \(2(x+y)=1+z\). 在高度 \(z\in[0,1]\) 处，三棱柱的截面为 \(x\geqslant0\)，\(y\geqslant0\)，\(x+y\leqslant1\) 的三角形；含 \(A\) 的一侧被截为 \(x+y\leqslant\frac{1+z}{2}\). 由于 \(\frac{1+z}{2}\in[\frac12,1]\)，该三角形与原截面相似，线性相似比为 \(\frac{1+z}{2}\)，面积比为 \(\left(\frac{1+z}{2}\right)^2\). 因而含 \(A\) 一侧的体积占总体积的 \(\int_0^1\left(\frac{1+z}{2}\right)^2\,\mathrm dz=\frac{1}{4}\int_0^1(1+2z+z^2)\,\mathrm dz=\frac{7}{12}\)，另一部分占 \(1-\frac{7}{12}=\frac{5}{12}\). 所以两部分体积之比为 \(7:5\).
\end{solution}

\section{解答题}

\begin{problem}
有四个数，其中前三个数成等差数列，后三个数成等比数列，并且第一个数与第四个数的和是 \(16\)，第二个数与第三个数的和是 \(12\). 求这四个数.
\end{problem}

\begin{solution}
设四个数依次为 \(a\)，\(b\)，\(c\)，\(d\). 由前三个数成等差数列，得 \(a+c=2b\)；由后三个数成等比数列，得 \(c^2=bd\). 再由题设的两个和式，得 \(a+d=16\)，\(b+c=12\). 于是 \(c=12-b\)，\(a=2b-c=3b-12\)，\(d=16-a=28-3b\). 代入 \(c^2=bd\)，得 \((12-b)^2=b(28-3b)\)，即 \(4b^2-52b+144=0\)，所以 \((b-4)(b-9)=0\). 当 \(b=4\) 时，\((a,c,d)=(0,8,16)\)；当 \(b=9\) 时，\((a,c,d)=(15,3,1)\). 这两组数分别满足等差、等比及两个和式条件，所以所求四个数为 \(0\)，\(4\)，\(8\)，\(16\)，或 \(15\)，\(9\)，\(3\)，\(1\).
\end{solution}

\begin{problem}
已知 \(\sin \alpha+\sin \beta=\frac{1}{4}\)，\(\cos \alpha+\cos \beta=\frac{1}{3}\)，求 \(\tan(\alpha+\beta)\) 的值.
\end{problem}

\begin{solution}
由和差化积公式，\(\sin\alpha+\sin\beta=2\sin\frac{\alpha+\beta}{2}\cos\frac{\alpha-\beta}{2}=\frac{1}{4}\)，\(\cos\alpha+\cos\beta=2\cos\frac{\alpha+\beta}{2}\cos\frac{\alpha-\beta}{2}=\frac{1}{3}\). 第二式右端不为零，所以 \(\cos\frac{\alpha-\beta}{2}\neq0\)，两式相除得 \(\tan\frac{\alpha+\beta}{2}=\frac{3}{4}\). 因而 \(\tan(\alpha+\beta)=\frac{2\tan\frac{\alpha+\beta}{2}}{1-\tan^2\frac{\alpha+\beta}{2}}=\frac{2\cdot\frac{3}{4}}{1-\left(\frac{3}{4}\right)^2}=\frac{24}{7}\).
\end{solution}

\begin{problem}
如图，在三棱锥 \(S-ABC\) 中，\(SA\perp\) 底面 \(ABC\)，\(AB\perp BC\). \(DE\) 垂直平分 \(SC\)，且分别交 \(AC\)，\(SC\) 于 \(D\)，\(E\). 又 \(SA=AB\)，\(SB=BC\). 求以 \(BD\) 为棱，以 \(BDE\) 与 \(BDC\) 为面的二面角的度数.

\bitmapfigure[max width=0.28\linewidth]{img/1990/national_science/q23_fig1.png}
\end{problem}

\begin{solution}
取长度单位使 \(SA=AB=1\). 由 \(SA\perp\) 底面 \(ABC\)，且 \(AB\perp BC\)，可取直角坐标系使 \(A=(0,0,0)\)，\(B=(1,0,0)\)，\(S=(0,0,1)\)，\(C=(1,\sqrt{2},0)\)，其中 \(SB=\sqrt{2}=BC\). 因为 \(DE\) 是 \(SC\) 的垂直平分线，故 \(E\) 是 \(SC\) 的中点，从而 \(E=\left(\frac{1}{2},\frac{\sqrt{2}}{2},\frac{1}{2}\right)\). 设 \(D=(t,t\sqrt{2},0)\) 在 \(AC\) 上，由 \(D\) 在 \(SC\) 的垂直平分线上，得 \(DS=DC\)，于是 \(DS^2=3t^2+1\)，\(DC^2=3(t-1)^2\)，所以 \(t=\frac{1}{3}\)，即 \(D=\left(\frac{1}{3},\frac{\sqrt{2}}{3},0\right)\).

因此 \(\overrightarrow{BD}=\left(-\frac{2}{3},\frac{\sqrt{2}}{3},0\right)\)，\(\overrightarrow{BE}=\left(-\frac{1}{2},\frac{\sqrt{2}}{2},\frac{1}{2}\right)\). 平面 \(BDE\) 的一个法向量为 \(\overrightarrow{BD}\times\overrightarrow{BE}=\left(\frac{\sqrt{2}}{6},\frac{1}{3},-\frac{\sqrt{2}}{6}\right)\)，而平面 \(BDC\) 就是底面 \(ABC\)，其法向量可取 \((0,0,1)\). 设所求二面角为 \(\theta\)，则 \(\cos\theta=\frac{\left| -\frac{\sqrt{2}}{6}\right|}{\sqrt{\frac{2}{36}+\frac{1}{9}+\frac{2}{36}}}=\frac{1}{2}\)，所以 \(\theta=60^\circ\).
\end{solution}

\begin{problem}
设 \(a\geqslant 0\)，在复数集 \(\mathbb C\) 中解方程：\(z^2+2|z|=a\).
\end{problem}

\begin{solution}
设 \(z=x+y\i\)，其中 \(x\)，\(y\in\R\)，记 \(r=|z|=\sqrt{x^2+y^2}\). 展开 \(z^2+2r=x^2-y^2+2r+2xy\i\)，与非负实数 \(a\) 比较虚部，得 \(2xy=0\)，所以必须分两种情况讨论.

当 \(y=0\) 时，\(z=x\)，方程化为 \(x^2+2|x|=a\). 令 \(t=|x|\geqslant0\)，则 \(t^2+2t=a\)，即 \((t+1)^2=a+1\). 因为 \(t+1>0\)，所以 \(t=\sqrt{a+1}-1\)，从而得到 \(z=\pm\left(\sqrt{a+1}-1\right)\).

当 \(x=0\) 时，\(z=y\i\)，方程化为 \(-y^2+2|y|=a\). 令 \(t=|y|\geqslant0\)，则 \(t^2-2t+a=0\)，所以 \(t=1\pm\sqrt{1-a}\). 这两个根为非负数且为实数，当且仅当 \(0\leqslant a\leqslant1\)，并给出 \(z=\pm\left(1+\sqrt{1-a}\right)\i\) 或 \(z=\pm\left(1-\sqrt{1-a}\right)\i\).

因此，\(a=0\) 时解为 \(z=0\)、\(z=\pm2\i\)；\(0<a<1\) 时解为 \(z=\pm\left(\sqrt{a+1}-1\right)\)、\(z=\pm\left(1+\sqrt{1-a}\right)\i\)、\(z=\pm\left(1-\sqrt{1-a}\right)\i\)；\(a=1\) 时解为 \(z=\pm\left(\sqrt{2}-1\right)\)、\(z=\pm\i\)；\(a>1\) 时解为 \(z=\pm\left(\sqrt{a+1}-1\right)\).
\end{solution}

\begin{problem}
设椭圆的中心是坐标原点，长轴在 \(x\) 轴上，离心率 \(e=\frac{\sqrt{3}}{2}\)，已知点 \(P\left( 0,\frac{3}{2} \right)\) 到这个椭圆上的点的最远距离是 \(\sqrt{7}\). 求这个椭圆的方程，并求椭圆上到点 \(P\) 的距离等于 \(\sqrt{7}\) 的点的坐标.
\end{problem}

\begin{solution}
设椭圆的半长轴、半短轴分别为 \(u\)，\(v\)，其中 \(u>v>0\). 由离心率条件，\(\frac{\sqrt{u^2-v^2}}{u}=\frac{\sqrt{3}}{2}\)，从而 \(v^2=\frac{u^2}{4}\)，即 \(u=2v\). 椭圆上的点可表示为 \((2v\cos t,v\sin t)\). 令 \(w=\sin t\)，则该点到 \(P\left(0,\frac{3}{2}\right)\) 的距离平方为 \(D^2=4v^2\cos^2t+\left(v\sin t-\frac{3}{2}\right)^2=4v^2-3v^2w^2-3vw+\frac{9}{4}\).

这是关于 \(w\) 的开口向下二次函数，且 \(-1\leqslant w\leqslant1\). 若 \(0<v<\frac{1}{2}\)，其最大值在 \(w=-1\) 处取得，且最大值为 \(\left(v+\frac{3}{2}\right)^2<4\)，不可能等于已知的 \(7\). 当 \(v=\frac{1}{2}\) 时最大值为 \(4\)，同样不符合条件，因此 \(v>\frac{1}{2}\). 此时二次函数的顶点 \(w=-\frac{1}{2v}\) 落在 \([-1,1]\) 内，最大距离平方为 \(4v^2+3\). 由最远距离为 \(\sqrt{7}\)，得 \(4v^2+3=7\)，所以 \(v=1\)，\(u=2\)，椭圆方程为 \(\frac{x^2}{4}+y^2=1\).

在该椭圆上，距离平方为 \(D^2=\frac{25}{4}-3w^2-3w\). 令其等于 \(7\)，得 \(3w^2+3w+\frac{3}{4}=0\)，即 \(\left(w+\frac{1}{2}\right)^2=0\). 所以 \(w=-\frac{1}{2}\)，从而 \(y=-\frac{1}{2}\)，\(x=\pm2\sqrt{1-w^2}=\pm\sqrt{3}\). 所求点的坐标为 \(\left(\pm\sqrt{3},-\frac{1}{2}\right)\).
\end{solution}

\begin{problem}
\(f(x)=\lg \frac{1+2^x+\cdots+(n-1)^x+n^xa}{n}\)，其中 \(a\) 是实数，\(n\) 是任意自然数且 \(n\geqslant 2\).
\begin{enumerate}
\item 如果 \(f(x)\) 当 \(x\in(-\infty,1]\) 时有意义，求 \(a\) 的取值范围；
\item 如果 \(a\in(0,1]\)，证明：\(2f(x)<f(2x)\) 当 \(x\neq 0\) 时成立.
\end{enumerate}
\end{problem}

\begin{solution}
\begin{enumerate}
\item 当 \(x\leqslant1\) 时，分子为 \(n^x\left(\sum_{k=1}^{n-1}\left(\frac{k}{n}\right)^x+a\right)\). 对每个 \(k=1\)，\(\ldots\)，\(n-1\)，都有 \(0<\frac{k}{n}<1\)，所以 \(\left(\frac{k}{n}\right)^x\) 随 \(x\) 增大而减小，从而在 \(x\leqslant1\) 时 \(\sum_{k=1}^{n-1}\left(\frac{k}{n}\right)^x\geqslant\sum_{k=1}^{n-1}\frac{k}{n}=\frac{n-1}{2}\). 因此，当且仅当 \(a+\frac{n-1}{2}>0\)，分子对所有 \(x\leqslant1\) 都为正，故 \(a>-\frac{n-1}{2}\).
\item 令 \(T(x)=1^x+2^x+\cdots+(n-1)^x+an^x\). 当 \(a\in(0,1]\) 时，所有 \(T(x)>0\)，且令 \(u_k=k^x\)（\(1\leqslant k\leqslant n-1\)），令 \(u_n=an^x\)，则 \(T(x)=\sum_{k=1}^{n}u_k\). 因为 \(0<a\leqslant1\)，有 \(T(2x)=\sum_{k=1}^{n-1}u_k^2+an^{2x}\geqslant\sum_{k=1}^{n}u_k^2\). 由柯西不等式，\(T(x)^2=\left(\sum_{k=1}^{n}u_k\right)^2\leqslant n\sum_{k=1}^{n}u_k^2\leqslant nT(2x)\). 当 \(0<a<1\) 时，第二个不等式严格；当 \(a=1\) 且 \(x\neq0\) 时，\(1^x\)，\(2^x\)，\(\ldots\)，\(n^x\) 不全相等，柯西不等式严格，所以对 \(x\neq0\) 总有 \(T(x)^2<nT(2x)\). 因而 \(f(2x)-2f(x)=\lg\frac{T(2x)}{n}-2\lg\frac{T(x)}{n}=\lg\frac{nT(2x)}{T(x)^2}>0\)，即 \(2f(x)<f(2x)\).
\end{enumerate}
\end{solution}
